\section{Response and information for the same detector}\label{sec:response}
\subsection{Fixed adaptive controllers}
A fixed controller means a parameter-independent Borel function of the
observed past, with its mode and gate commands known to the controller.
When differentiating the unknown physical radius, this function is held
fixed. Its thresholds need not be differentiable as functions of the past.
There is no conditioning on a final adaptively chosen design as though it
had been fixed in advance.

For $0\le j\le3$ set
\[
 \beta_j=\sup_{R,a,x}\frac{|\partial_R^jk_R^a(x)|}{k_R^a(x)},
 \qquad \beta_0=1,
\]
and put $\beta_j=0$ for $j>3$. These constants are finite by
Proposition~\ref{prop:raw}. They are explicit: differentiate the three
polynomials \eqref{eq:raw-kernel} and divide a coefficient upper bound by
$h_*$. No derivative of a grazing collision time occurs in this estimate.
For an accepted density, the factor $g(x)$ cancels in its derivative ratio.
For failure, positivity gives
\[
 |\partial_R^jf_R^{a,g}|
 \le\int (1-g)|\partial_R^jk_R^a|\,d\nu\le\beta_j f_R^{a,g}.
\]

\begin{theorem}[Polynomial path response for every fixed controller]\label{thm:response}
For a length-$N$ fixed adaptive controller, the density of its readout
history with respect to $\mu^N$ is a polynomial in $R$ of degree at most
$3N$. Its path law $P_R$ is strongly $C^r$ in signed-measure variation for
every finite $r$. With the variation norm $\|\sigma\|_{\rm var}=\int|d\sigma|$,
\begin{equation}
 \norm{\partial_R^r P_R}_{\rm var}
 \le r!\sum_{j_1+\cdots+j_N=r}\prod_{i=1}^N\frac{\beta_{j_i}}{j_i!}.
 \label{eq:response-bound}
\end{equation}
The derivatives above order $3N$ vanish. At interior radii the family is
differentiable in quadratic mean. Its score is the sum of the conditional
one-cartridge scores and its Fisher information is
\begin{equation}
 I_N(R)=\sum_{i=1}^N\E_R\bigl[I(R;a_{H_{i-1}},g_{H_{i-1}})\bigr].
 \label{eq:adaptive-info}
\end{equation}
All statements include both placement failure and censoring.
\end{theorem}
\begin{proof}
At a fixed realized report history, each command is a fixed value, the same
for every possible radius. The chronological density is the product of the
corresponding polynomial conditional densities. Coefficients of these
polynomials are measurable in the history and uniformly bounded over all
commands, since the modes are compact and the gates bounded. The reference
measure $\mu^N$ has finite mass $4^N$. The density is thus a polynomial with
$L^1$ coefficient functions, proving strong differentiability and exact
vanishing at high order. Applying the product rule pointwise and dividing
by the positive path density bounds each term by the product of its
$\beta$ ratios. Integration proves \eqref{eq:response-bound}.

The conditional densities, and hence the path densities, have common
support independent of the radius, even for gates taking values zero or one
as in Theorem~\ref{thm:boundary}. On that support their likelihood ratios
between any two radii are uniformly bounded. Their first two relative
derivatives are bounded by the product-rule estimates above. Thus
$\partial_R^2\sqrt{p_R}=\sqrt{p_R}
\{p_R''/(2p_R)-(p_R'/p_R)^2/4\}$ is bounded in absolute value by a
constant times $\sqrt{p_{R_0}}$ for a fixed reference radius $R_0$.
Off the common support all derivatives are zero. Taylor's integral
remainder then gives an $L^2(\mu^N)$ error $O(h^2)$ after the first-order
square-root term. This proves quadratic-mean differentiability with score
$\partial_R\log p_R$ on the common support. Differentiating the logarithm of the chronological
product gives the sum of conditional scores. Each has conditional mean
zero, because its positive conditional density integrates to one and has
an integrable derivative. Earlier scores are measurable with respect to
the preceding history, so distinct increments are orthogonal in $L^2$.
Taking their squared sum proves \eqref{eq:adaptive-info}.
\end{proof}

A policy optimized for one nominal specification can be held fixed and
analyzed by this theorem. A pointwise maximum of policy values need not be
differentiable when maximizing policies exchange. We prove continuity and
attainment of that maximum in Section~\ref{sec:control}; we do not apply
the fixed-controller derivative to a changing optimizer.

\subsection{A strictly informative collision-coordinate mark}
The raw cartridge report has a placement flag with probability
$s_R=\pi R^2/a_0$. Conditional on placement success it has the equilibrium
collision probability $p_R=2TR/A(R)$. Conditional on a collision the mark
density with respect to $dy/(2\pi)$ is $1+c_R\cos(y-\theta)$, where
$c_R=\epsilon R^2$. For $|c|<1$,
\[
 \int\frac{dy}{2\pi(1+c\cos y)}=\frac1{\sqrt{1-c^2}}.
\]
For example the substitution $t=\tan(y/2)$, splitting at $y=\pi$, reduces
the integral to the integral of
$\{(1+c)+(1-c)t^2\}^{-1}$ over the real line, divided by $\pi$.
The identity
$c^2\cos^2y/(1+c\cos y)=c\cos y-1+(1+c\cos y)^{-1}$ then gives the
conditional mark information
\begin{equation}
 J_R=\frac4{R^2}\left(\frac1{\sqrt{1-\epsilon^2R^4}}-1\right),
 \qquad J_R=0\text{ when }\epsilon=0.
 \label{eq:mark-info}
\end{equation}
The two conditional normalization identities center the corresponding
scores. Thus direct differentiation, or orthogonality of the conditional
scores, yields
\begin{equation}
 I_{\rm raw}(R)
 =\frac{(s_R')^2}{s_R(1-s_R)}
 +(1-s_R)\left\{\frac{(p_R')^2}{p_R(1-p_R)}+p_R J_R\right\}.
 \label{eq:raw-info}
\end{equation}
In particular, when $\epsilon>0$ the mark adds strictly positive information
over observing the same placement and collision flags without the mark.
The comparison uses the same counted experiment; it is not a comparison
between an uncharged equilibrium trial and an ambient placement attempt.

For a gated report, the accepted score equals the raw score, whereas the
censoring score is $f_R'/f_R$ whenever $f_R>0$. If $f_R=0$, the
common-null censoring branch is omitted and its information contribution is
zero. More explicitly, on the positive-failure branch,
\[
 \frac{f_R'}{f_R}=\frac{\int(1-g)k_R^a\,\partial_R\log k_R^a\,d\nu}
                       {\int(1-g)k_R^a\,d\nu}.
\]
It is the conditional mean of the raw score given censoring. Jensen's
inequality, or expanding the conditional variance, proves
$I(R;a,g)\le I_{\rm raw}(R;a)$. The information on the failure atom has not
been set to zero. Since their likelihood ratios on the common support are uniformly bounded
above and below, their relative entropies are finite
for all radii in $I$. This conclusion holds for this detector, in contrast
to the ideal geometric record analyzed in Section~\ref{sec:lab}.

\subsection{Sources and observables}
For bounded parameter-independent path tests $F,V$ and a fixed controller,
$P_R(e^VF)$ and $P_R(e^V)$ are polynomials of degree at most $3N$ with
explicit derivative integrals. The normalized tilted expectation is their
ratio, with a positive denominator. At order $r$ its derivatives are given
recursively from $ZU=J$ by
\[
 U^{(r)}=Z^{-1}\left[J^{(r)}-
          \sum_{j=1}^r\binom rj Z^{(j)}U^{(r-j)}\right].
\]
If the tests also depend smoothly on $R$ in supremum norm, retain their
ordinary product derivatives. These formulas apply to the actual readout
law; they do not manufacture a score for a moving singular mechanical
record. A measurable threshold of the noisy readout remains an admissible
bounded test, even when it is not a smooth test of the ideal collision path.
